\documentclass[11pt,a4paper]{article}

% ------------------------------------------------------------------
%  Modelos Avanzados de Computacion - Relacion 4 (VERSION COMENTADA)
%  Resolucion con fines de estudio + explicacion de cada concepto.
% ------------------------------------------------------------------

\usepackage[utf8]{inputenc}
\usepackage[T1]{fontenc}
\usepackage{babel}
\babelprovide[import, main]{spanish}
\usepackage{amsmath,amssymb,amsthm}
\usepackage{mathtools}
\usepackage{enumitem}
\usepackage{geometry}
\usepackage{xcolor}
\usepackage{tcolorbox}
\tcbuselibrary{skins,breakable}
\usepackage{tikz}
\usetikzlibrary{arrows.meta,positioning,automata}
\usepackage{fancyhdr}
\usepackage{hyperref}
\hypersetup{colorlinks=true,linkcolor=blue,urlcolor=blue}

\geometry{margin=2.4cm}

% --- Entornos de teoremas -----------------------------------------
\theoremstyle{plain}
\newtheorem{teorema}{Teorema}
\newtheorem{lema}{Lema}
\newtheorem{proposicion}{Proposici\'on}
\theoremstyle{definition}
\newtheorem{definicion}{Definici\'on}
\newtheorem*{nota}{Nota}
\newtheorem*{idea}{Idea}

% --- Cajas pedagogicas --------------------------------------------
% Caja "concepto": explica un termino o herramienta de fondo.
\newtcolorbox{concepto}[1][]{
  breakable, enhanced, colback=blue!4!white, colframe=blue!55!black,
  fonttitle=\bfseries, title={$\triangleright$~Concepto: #1},
  boxrule=0.6pt, left=3mm, right=3mm, top=1.5mm, bottom=1.5mm,
  before skip=3mm, after skip=3mm}

% Caja "porque": explica el motivo de un paso concreto.
\newtcolorbox{porque}{
  breakable, enhanced, colback=green!4!white, colframe=green!45!black,
  fonttitle=\bfseries, title={\textbf{\textquestiondown{}Por qu\'e este paso?}},
  boxrule=0.6pt, left=3mm, right=3mm, top=1.5mm, bottom=1.5mm,
  before skip=3mm, after skip=3mm}

% Caja "cuidado": errores tipicos / matices delicados.
\newtcolorbox{cuidado}{
  breakable, enhanced, colback=red!4!white, colframe=red!55!black,
  fonttitle=\bfseries, title={\textbf{\textexclamdown{}Cuidado!}},
  boxrule=0.6pt, left=3mm, right=3mm, top=1.5mm, bottom=1.5mm,
  before skip=3mm, after skip=3mm}

% --- Macros utiles ------------------------------------------------
\newcommand{\NP}{\mathrm{NP}}
\newcommand{\coNP}{\mathrm{CoNP}}
\newcommand{\Pclass}{\mathrm{P}}
\newcommand{\Lclass}{\mathrm{L}}
\newcommand{\SPACE}{\mathrm{ESPACIO}}
\newcommand{\TIME}{\mathrm{TIEMPO}}
\newcommand{\reduce}{\propto}
\newcommand{\Lstar}{L^{*}}
\renewcommand{\qedsymbol}{$\blacksquare$}

% --- Encabezados --------------------------------------------------
\pagestyle{fancy}
\fancyhf{}
\lhead{\small Modelos Avanzados de Computaci\'on}
\rhead{\small Relaci\'on 4 (comentada)}
\cfoot{\thepage}

\title{\textbf{Modelos Avanzados de Computaci\'on}\\[2mm]
\large Relaci\'on 4 --- Resoluci\'on comentada para estudio}
\author{}
\date{}

\begin{document}
\maketitle
\thispagestyle{fancy}

\begin{nota}
Esta es la versi\'on \textbf{comentada} de la relaci\'on. Adem\'as de las
demostraciones, cada ejercicio incluye:
\begin{itemize}[nosep]
  \item cajas \textcolor{blue!55!black}{\textbf{Concepto}} que explican desde cero
        las herramientas que se usan (m\'aquinas de Turing, espacio logar\'itmico,
        certificados, reducciones, etc.), sin dar nada por sabido;
  \item cajas \textcolor{green!45!black}{\textbf{\textquestiondown{}Por qu\'e este paso?}} que
        justifican \emph{por qu\'e} se hace cada cosa de esa manera;
  \item cajas \textcolor{red!55!black}{\textbf{\textexclamdown{}Cuidado!}} con los errores t\'ipicos
        y los matices delicados que suelen caer en examen.
\end{itemize}
Los resultados son est\'andar en teor\'ia de la complejidad (Sipser, Papadimitriou,
Garey--Johnson); donde un argumento sigue de cerca esos textos se indica.
\end{nota}

\tableofcontents
\vspace{4mm}\hrule\vspace{4mm}

% ==================================================================
\section*{Conceptos base (l\'eelo antes de empezar)}
\addcontentsline{toc}{section}{Conceptos base}
% ==================================================================

Casi toda la relaci\'on gira en torno a unas pocas ideas. Las reunimos aqu\'i para no
repetirlas en cada ejercicio.

\begin{concepto}[M\'aquina de Turing (MT) y por qu\'e medimos ``recursos'']
Una \textbf{m\'aquina de Turing} es el modelo matem\'atico de ``ordenador ideal'':
una o varias cintas infinitas divididas en casillas, una cabeza que lee/escribe
s\'imbolos y se mueve, y un control finito de estados. Todo algoritmo se puede
describir como una MT. La usamos porque permite \emph{medir con precisi\'on} dos
recursos:
\begin{itemize}[nosep]
  \item \textbf{Tiempo:} n\'umero de pasos hasta parar.
  \item \textbf{Espacio (memoria):} n\'umero de casillas distintas que se llegan a
        usar.
\end{itemize}
Cuando decimos ``tiempo polin\'omico'' queremos decir: el n\'umero de pasos est\'a
acotado por $n^{c}$ para alguna constante $c$, donde $n$ es la longitud de la
entrada. ``Espacio logar\'itmico'' significa que la memoria usada es $O(\log n)$.
\end{concepto}

\begin{concepto}[Cinta de entrada de \emph{solo lectura} (clave para el espacio)]
Para medir espacios \emph{sublineales} (menos que $n$) se usa un modelo con
\textbf{dos cintas}:
\begin{itemize}[nosep]
  \item una \textbf{cinta de entrada de solo lectura}: contiene $x$, la cabeza puede
        moverse libremente por ella y releer, \emph{pero no se cuenta} como memoria
        (es ``de fuera'');
  \item una \textbf{cinta de trabajo} de lectura/escritura: \emph{esta s\'i} se
        cuenta, y es donde medimos el espacio.
\end{itemize}
\textquestiondown{}Por qu\'e esta separaci\'on? Porque si contaramos la cinta de entrada, ning\'un
algoritmo podr\'ia usar menos de $n$ espacio (solo guardar la entrada ya cuesta $n$).
Al no contarla, tiene sentido hablar de algoritmos que usan $O(\log n)$ memoria
\emph{de trabajo} aunque la entrada sea grande. Esto es lo que hace posible la clase
$\Lclass$.
\end{concepto}

\begin{concepto}[Por qu\'e un \'indice cuesta $O(\log n)$ bits]
Una idea que se repite muchas veces: para guardar un n\'umero (un contador, una
posici\'on, un \'indice) entre $0$ y $n$, se necesitan $\lceil\log_2(n+1)\rceil$
bits en binario. Por eso ``llevar unos cuantos punteros/contadores acotados por $n$''
ocupa $O(\log n)$ espacio. Es el truco central de los ejercicios 5, 10 y 11.
\end{concepto}

\begin{concepto}[Las clases $\Pclass$, $\NP$, $\coNP$, $\Lclass$, $\SPACE(n)$]
\begin{itemize}[nosep]
  \item $\Pclass$ = problemas decidibles por una MT \emph{determinista} en
        \textbf{tiempo} polin\'omico. Intuici\'on: ``resolubles de forma eficiente''.
  \item $\NP$ = problemas en los que, si la respuesta es ``S\'i'', existe una
        \textbf{prueba corta} (un \emph{certificado}) que se puede \textbf{verificar}
        en tiempo polin\'omico. Intuici\'on: ``f\'acil de comprobar, quiz\'a dif\'icil
        de encontrar''.
  \item $\coNP$ = los \emph{complementarios} de los problemas de $\NP$ (intercambiar
        ``S\'i'' y ``No''). Aqu\'i lo f\'acil de certificar es el ``No''.
  \item $\Lclass = \SPACE(\log n)$ = problemas decidibles en \textbf{espacio}
        logar\'itmico (de trabajo).
  \item $\SPACE(n)$ = problemas decidibles en \textbf{espacio} lineal (memoria de
        trabajo $O(n)$).
\end{itemize}
Relaciones b\'asicas: $\Lclass\subseteq\Pclass\subseteq\NP$ y
$\Lclass\subseteq\SPACE(n)$. Cu\'ales de estas inclusiones son estrictas es, en
general, un problema abierto (incluido el famoso $\Pclass$ vs $\NP$).
\end{concepto}

\begin{concepto}[Certificado y verificador (la definici\'on operativa de $\NP$)]
Un problema $P$ est\'a en $\NP$ si existe un \textbf{verificador} $V$ (algoritmo
determinista polin\'omico) y un polinomio $q$ tales que:
\[
P(x)=\text{``S\'i''} \iff \exists\, c \text{ con } |c|\le q(|x|)\ \text{tal que }
V(x,c)=1.
\]
La cadena $c$ es el \textbf{certificado} (o ``testigo''): una pista que demuestra que
la respuesta es ``S\'i''. La condici\'on $|c|\le q(|x|)$ obliga a que la pista sea
\emph{corta} (polin\'omica). Ejemplo: para ``\textquestiondown{}es este grafo coloreable con 3
colores?'', un certificado es la propia asignaci\'on de colores; verificarla
(mirar que ninguna arista tenga dos extremos del mismo color) es trivial y
r\'apido. Esta forma de ver $\NP$ se usa en los ejercicios 7 y 8.
\end{concepto}

\begin{concepto}[Reducci\'on en espacio logar\'itmico ($P_1\reduce P_2$)]
Reducir $P_1$ a $P_2$ significa \emph{transformar} cualquier entrada $x$ de $P_1$ en
una entrada $f(x)$ de $P_2$ de modo que la respuesta no cambie:
\[
P_1(x)=\text{``S\'i''}\iff P_2(f(x))=\text{``S\'i''}.
\]
En esta asignatura se exige adem\'as que $f$ se pueda calcular en \textbf{espacio
logar\'itmico}. Intuici\'on: ``si s\'e resolver $P_2$, y puedo traducir baratamente
$P_1$ a $P_2$, entonces tambi\'en s\'e resolver $P_1$''. Es la herramienta para
comparar la dificultad de problemas y la base del ejercicio 13.
\end{concepto}

\begin{concepto}[Notaci\'on $O(\cdot)$ y ``$=o(\cdot)$'']
$g(n)=O(h(n))$ significa que, a partir de cierto punto, $g$ no crece m\'as r\'apido
que $h$ salvo un factor constante (``$g$ est\'a dominada por $h$''). $g(n)=o(h(n))$ es
m\'as fuerte: $g/h\to 0$, es decir $h$ crece \emph{estrictamente} m\'as r\'apido.
Estas notaciones se usan en todo el ejercicio 4 y en el teorema de jerarqu\'ia del 13.
\end{concepto}

\vspace{2mm}\hrule\vspace{4mm}

% ==================================================================
\section{Grafos dirigidos ac\'iclicos (DAG)}
% ==================================================================

\begin{concepto}[Vocabulario de grafos dirigidos]
Un \textbf{grafo dirigido} $G=(V,E)$ tiene nodos (v\'ertices) $V$ y arcos
\emph{con sentido} $E$: un arco $(u,v)$ va ``de $u$ hacia $v$''. Un \textbf{ciclo
dirigido} es un camino que respeta los sentidos y vuelve al punto de partida
($v_1\to v_2\to\cdots\to v_k\to v_1$). \textbf{Ac\'iclico} = sin ciclos dirigidos
(se abrevia DAG, \emph{directed acyclic graph}). Una \textbf{fuente} es un nodo
sin arcos \emph{entrantes} (nadie apunta a \'el); un \textbf{sumidero} es uno sin
arcos \emph{salientes}.
\end{concepto}

\subsection*{(a) Todo DAG finito tiene una fuente}

\begin{proposicion}
Todo grafo dirigido finito y ac\'iclico posee al menos una \emph{fuente}.
\end{proposicion}

\begin{porque}
Queremos una fuente porque ser\'a el ``primer'' nodo del orden que construiremos en
(b). La estrategia es t\'ipica: suponemos que \emph{no} existe y derivamos un
absurdo (un ciclo), contradiciendo que el grafo sea ac\'iclico.
\end{porque}

\begin{proof}
Razonamos por contradicci\'on. Sea $G=(V,E)$ finito y ac\'iclico, y supongamos que
\emph{ning\'un} nodo es fuente; es decir, todo nodo tiene al menos un arco entrante.

Construimos un camino hacia atr\'as. Tomamos un nodo cualquiera $v_0$. Como $v_0$ no
es fuente, existe un arco $(v_1,v_0)\in E$. Como $v_1$ tampoco es fuente, existe
$(v_2,v_1)\in E$. Repitiendo, generamos una sucesi\'on
\[
\cdots \to v_3 \to v_2 \to v_1 \to v_0,
\]
en la que cada $v_{i+1}$ tiene un arco hacia $v_i$. Como $|V|=n$ es finito, entre
los $n+1$ primeros nodos $v_0,v_1,\dots,v_n$ debe haber, por el principio del
palomar, dos \'indices $i<j$ con $v_i=v_j$. El tramo
\[
v_j \to v_{j-1}\to \cdots \to v_i = v_j
\]
es entonces un \textbf{ciclo dirigido}, lo que contradice que $G$ sea ac\'iclico.
Por tanto alg\'un nodo no tiene arco entrante: $G$ tiene una fuente.
\end{proof}

\begin{concepto}[Principio del palomar]
Si repartimos $n+1$ objetos en $n$ cajas, alguna caja recibe $\ge 2$ objetos. Aqu\'i:
tenemos $n+1$ nodos en la sucesi\'on, pero solo $n$ nodos distintos disponibles en el
grafo, as\'i que dos de ellos \emph{tienen} que coincidir. Es lo que cierra el ciclo.
\end{concepto}

\begin{nota}
Por simetr\'ia (invirtiendo todos los arcos) el mismo argumento prueba que todo DAG
finito tiene tambi\'en un \emph{sumidero}.
\end{nota}

\subsection*{(b) Numeraci\'on topol\'ogica}

\begin{proposicion}
Un grafo dirigido con $n$ nodos es ac\'iclico si y solo si se pueden numerar sus nodos
$1,2,\dots,n$ de modo que todo arco vaya de un n\'umero menor a uno mayor.
\end{proposicion}

\begin{concepto}[Demostrar un ``si y solo si'' ($\Leftrightarrow$)]
Una equivalencia ``$A\Leftrightarrow B$'' se prueba en \textbf{dos direcciones}
separadas: ``$A\Rightarrow B$'' y ``$B\Rightarrow A$''. Aqu\'i $A=$``es ac\'iclico'' y
$B=$``existe la numeraci\'on''. Hacemos primero la direcci\'on f\'acil
($B\Rightarrow A$) y luego la constructiva ($A\Rightarrow B$).
\end{concepto}

\begin{proof}
$(\Leftarrow)$ Supongamos que existe tal numeraci\'on $\nu:V\to\{1,\dots,n\}$ con
$\nu(u)<\nu(v)$ para todo arco $(u,v)$. Si hubiera un ciclo
$w_1\to w_2\to\cdots\to w_k\to w_1$, encadenando las desigualdades obtendr\'iamos
\[
\nu(w_1)<\nu(w_2)<\cdots<\nu(w_k)<\nu(w_1),
\]
es decir $\nu(w_1)<\nu(w_1)$, absurdo. Luego $G$ es ac\'iclico.

$(\Rightarrow)$ Supongamos $G$ ac\'iclico. Construimos la numeraci\'on por inducci\'on
sobre $n$.

Por el apartado (a), $G$ tiene una fuente $s$. Le asignamos el n\'umero $1$:
$\nu(s)=1$. Consideramos $G'=G\setminus\{s\}$ (eliminamos $s$ y sus arcos
salientes). $G'$ sigue siendo ac\'iclico y tiene $n-1$ nodos; por hip\'otesis de
inducci\'on admite una numeraci\'on $\nu':V\setminus\{s\}\to\{1,\dots,n-1\}$
compatible. Definimos $\nu(v)=\nu'(v)+1$ para $v\neq s$.

Comprobamos que todo arco va de menor a mayor:
\begin{itemize}[nosep]
  \item Arcos que salen de $s$: van de $\nu(s)=1$ a alg\'un $\nu(v)\geq 2$. Correcto.
  \item Arcos que no tocan $s$: estaban en $G'$ y eran crecientes para $\nu'$, luego
        siguen siendo crecientes para $\nu=\nu'+1$. Correcto.
  \item No hay arcos \emph{entrantes} a $s$ porque $s$ es fuente.
\end{itemize}
Esto completa la inducci\'on.
\end{proof}

\begin{concepto}[Inducci\'on sobre el tama\~no, y qu\'e es un ``orden topol\'ogico'']
La \textbf{inducci\'on} funciona aqu\'i porque al quitar la fuente obtenemos un grafo
m\'as peque\~no \emph{del mismo tipo} (sigue siendo DAG): resolvemos el caso peque\~no
y ``pegamos'' la fuente delante. El orden resultante se llama \textbf{orden
topol\'ogico}: es exactamente ``colocar los nodos en fila de forma que todas las
flechas apunten hacia la derecha''. Existe \emph{si y solo si} el grafo es ac\'iclico,
que es justo lo que dice esta proposici\'on.
\end{concepto}

\begin{porque}
\textquestiondown{}Por qu\'e $G'$ sigue siendo ac\'iclico? Porque \textbf{quitar} nodos o arcos nunca
puede \emph{crear} un ciclo (un ciclo en $G'$ ya estar\'ia en $G$). Esto es lo que
nos permite aplicar la hip\'otesis de inducci\'on sin problemas.
\end{porque}

\subsection*{(c) Algoritmo polin\'omico para decidir si un grafo es ac\'iclico}

La prueba de (b) es constructiva y da directamente el algoritmo (eliminaci\'on
iterativa de fuentes, conocido como algoritmo de Kahn):

\begin{quote}
\textbf{Algoritmo} \textsc{EsAciclico}$(G=(V,E))$:
\begin{enumerate}[nosep]
  \item Calcular el grado de entrada $d^{-}(v)$ de cada nodo. Inicializar una cola
        $Q$ con todos los nodos de grado de entrada $0$ (las fuentes). Contador
        $c\leftarrow 0$.
  \item Mientras $Q\neq\emptyset$:
        \begin{enumerate}[nosep]
          \item Extraer $u$ de $Q$; $c\leftarrow c+1$.
          \item Para cada arco $(u,w)$: decrementar $d^{-}(w)$; si queda en $0$,
                a\~nadir $w$ a $Q$.
        \end{enumerate}
  \item Si $c=|V|$ devolver \textsc{Aciclico}; en otro caso devolver
        \textsc{Tiene Ciclo}.
\end{enumerate}
\end{quote}

\begin{concepto}[``Grado de entrada'' y por qu\'e lo decrementamos]
El \textbf{grado de entrada} $d^{-}(v)$ es el n\'umero de arcos que apuntan a $v$. Un
nodo es fuente \emph{ahora mismo} si $d^{-}(v)=0$. Cuando ``eliminamos'' un nodo $u$,
sus arcos salientes desaparecen, as\'i que cada vecino $w$ pierde un arco entrante:
por eso hacemos $d^{-}(w)\!\mathrel{-}\!=\!1$. Si con eso $w$ se queda a $0$, $w$ se
ha convertido en una nueva fuente y entra en la cola. As\'i simulamos ``quitar la
fuente y repetir'' sin tocar el grafo f\'isicamente.
\end{concepto}

\paragraph{Correcci\'on.} Si $G$ es ac\'iclico, por (a) en cada paso del subgrafo
restante siempre queda una fuente, as\'i que el proceso elimina los $n$ nodos y
$c=n$. Rec\'iprocamente, si el proceso se detiene con $c<n$, todos los nodos
restantes tienen grado de entrada $\geq 1$, luego (por el argumento de (a) aplicado
al subgrafo restante) ese subgrafo contiene un ciclo, y por tanto $G$ tambi\'en.

\paragraph{Complejidad.} Cada arco se procesa exactamente una vez (al eliminar su
origen) y cada nodo se encola/desencola una vez. El coste es $O(|V|+|E|)$, lineal y
en particular polin\'omico.

\begin{porque}
\textquestiondown{}Por qu\'e nos vale ``polin\'omico'' y no buscamos algo m\'as r\'apido? Porque el
enunciado solo pide \emph{demostrar que el problema es tratable} (est\'a en
$\Pclass$). $O(|V|+|E|)$ es de hecho \'optimo (hay que mirar toda la entrada), pero
lo \'unico que el ejercicio exige es la cota polin\'omica.
\end{porque}

% ==================================================================
\section{Grafos bipartitos}
% ==================================================================

\begin{concepto}[Grafo no dirigido, ciclo, longitud, bipartito]
Ahora los arcos no tienen sentido: una \textbf{arista} $\{u,v\}$ une dos v\'ertices.
La \textbf{longitud} de un ciclo es su n\'umero de aristas. Un grafo es
\textbf{bipartito} si sus v\'ertices se pueden separar en dos grupos $A$ y $B$
(no necesariamente del mismo tama\~no) de modo que \emph{toda} arista tenga un
extremo en $A$ y el otro en $B$ --- nunca dos del mismo grupo. Pensar en
``2-coloreado'': pintar cada v\'ertice de rojo o azul sin que dos vecinos compartan
color.
\end{concepto}

\subsection*{(a) Bipartito $\Leftrightarrow$ todos los ciclos de longitud par}

\begin{proposicion}
Un grafo $G=(V,E)$ es bipartito si y solo si todos sus ciclos tienen longitud par.
\end{proposicion}

\begin{proof}
$(\Rightarrow)$ Supongamos $G$ bipartito con partici\'on $V=A\,\dot\cup\,B$. Sea
$v_0,v_1,\dots,v_k=v_0$ un ciclo. Cada arista cambia de lado: si $v_0\in A$ entonces
$v_1\in B$, $v_2\in A$, \dots; en general $v_i\in A$ si $i$ es par y $v_i\in B$ si $i$
es impar. Como el ciclo se cierra en $v_k=v_0\in A$, el \'indice $k$ ha de ser par.
Luego todo ciclo tiene longitud par.

$(\Leftarrow)$ Supongamos que todos los ciclos son de longitud par. Trabajamos por
componentes conexas (la uni\'on de las biparticiones de cada componente da una
biparticion del grafo entero). Fijada una componente, elegimos una ra\'iz $r$ y
definimos, para cada v\'ertice $v$, la distancia $d(v)=$ longitud del camino m\'as
corto de $r$ a $v$. Coloreamos:
\[
A=\{v: d(v)\text{ par}\},\qquad B=\{v: d(v)\text{ impar}\}.
\]
Veamos que ninguna arista une dos v\'ertices del mismo color. Sea $\{u,v\}\in E$. Si
fuera $d(u)\equiv d(v)\pmod 2$, considere\-mos los caminos m\'as cortos $P_u$ de $r$ a
$u$ y $P_v$ de $r$ a $v$. Sea $w$ su \'ultimo v\'ertice com\'un. Los tramos
$w\!\to\!u$ y $w\!\to\!v$ tienen longitudes $d(u)-d(w)$ y $d(v)-d(w)$, de la
\emph{misma} paridad. El recorrido
\[
w\to\cdots\to u \;-\; v\to\cdots\to w
\]
(bajar a $u$, usar la arista $\{u,v\}$, subir desde $v$) es un ciclo de longitud
$\big(d(u)-d(w)\big) + 1 + \big(d(v)-d(w)\big)$, que es \textbf{impar} (suma de dos
n\'umeros de igual paridad m\'as $1$). Esto contradice la hip\'otesis. Por tanto toda
arista une colores distintos y la componente es bipartita.
\end{proof}

\begin{concepto}[``$\pmod 2$'' y paridad]
``$d(u)\equiv d(v)\pmod 2$'' significa que $d(u)$ y $d(v)$ tienen la \emph{misma
paridad} (ambos pares o ambos impares). La diferencia de dos n\'umeros de igual
paridad es par; al sumarle $1$ queda impar. Esa es toda la aritm\'etica que usa la
prueba.
\end{concepto}

\begin{porque}
\textbf{La idea profunda:} en un grafo bipartito, la paridad de la distancia a la
ra\'iz determina el color, y \emph{nunca} cambia de forma inconsistente. Un ciclo
impar es justo el obst\'aculo que har\'ia imposible asignar colores de forma
coherente (al recorrerlo volver\'iamos al inicio con el color cambiado). Por eso
``sin ciclos impares'' $\equiv$ ``coloreable con 2 colores'' $\equiv$ ``bipartito''.
\end{porque}

\subsection*{(b) Algoritmo polin\'omico para comprobar bipartici\'on}

\begin{concepto}[BFS, recorrido en anchura]
\textbf{BFS} (\emph{breadth-first search}) explora el grafo ``por capas'': primero la
ra\'iz, luego sus vecinos, luego los vecinos de estos, etc. Garantiza que cada
v\'ertice se alcanza por su camino m\'as corto, de ah\'i que conozcamos la paridad de
$d(v)$ de forma natural mientras recorremos. Coste $O(|V|+|E|)$.
\end{concepto}

\begin{quote}
\textbf{Algoritmo} \textsc{EsBipartito}$(G)$:
\begin{enumerate}[nosep]
  \item Marcar todos los v\'ertices como sin color.
  \item Para cada v\'ertice $r$ a\'un sin color (cabecera de una componente):
        \begin{enumerate}[nosep]
          \item Asignar $color(r)=0$ y lanzar un BFS desde $r$.
          \item Al visitar una arista $\{u,v\}$: si $v$ no tiene color, asignar
                $color(v)=1-color(u)$; si ya lo tiene y $color(v)=color(u)$,
                devolver \textsc{No bipartito}.
        \end{enumerate}
  \item Si nunca hubo conflicto, devolver \textsc{Bipartito} con la partici\'on dada
        por los colores.
\end{enumerate}
\end{quote}

\paragraph{Correcci\'on.} El BFS asigna a cada v\'ertice la paridad de su distancia a
la ra\'iz, exactamente el coloreado de (a). Un conflicto ($color(u)=color(v)$ en una
arista) corresponde a la existencia de un ciclo impar; su ausencia certifica una
biparticion v\'alida.

\paragraph{Complejidad.} Recorrido BFS est\'andar: $O(|V|+|E|)$, polin\'omico.

% ==================================================================
\section{$\Pclass$ es cerrada para uni\'on e intersecci\'on}
% ==================================================================

\begin{concepto}[Qu\'e significa ``cerrada para una operaci\'on'']
Decir que una clase $\mathcal{C}$ es \textbf{cerrada} para una operaci\'on significa
que, al aplicar esa operaci\'on a elementos de $\mathcal{C}$, el resultado
\emph{sigue} en $\mathcal{C}$. Aqu\'i: si $L_1,L_2\in\Pclass$, queremos que
$L_1\cup L_2$ y $L_1\cap L_2$ tambi\'en est\'en en $\Pclass$. Es como comprobar que
``los pares son cerrados para la suma'': par $+$ par $=$ par.
\end{concepto}

\begin{teorema}
Si $L_1,L_2\in\Pclass$, entonces $L_1\cup L_2\in\Pclass$ y
$L_1\cap L_2\in\Pclass$.
\end{teorema}

\begin{proof}
Por hip\'otesis existen m\'aquinas deterministas $M_1$ y $M_2$ que deciden $L_1$ y
$L_2$ en tiempo $O(n^{c_1})$ y $O(n^{c_2})$. Construimos $M$ que, con entrada $x$ de
longitud $n$:
\begin{enumerate}[nosep]
  \item Simula $M_1(x)$ y guarda su respuesta $b_1\in\{\text{S\'i},\text{No}\}$.
  \item Simula $M_2(x)$ y guarda su respuesta $b_2$.
  \item \textbf{Uni\'on:} acepta si $b_1=\text{S\'i}$ \emph{o} $b_2=\text{S\'i}$.\\
        \textbf{Intersecci\'on:} acepta si $b_1=\text{S\'i}$ \emph{y}
        $b_2=\text{S\'i}$.
\end{enumerate}
$M$ es determinista y siempre para. Su tiempo es a lo sumo
$O(n^{c_1})+O(n^{c_2})+O(1)=O(n^{\max(c_1,c_2)})$, polin\'omico. Y
$L(M)=L_1\cup L_2$ (resp. $L_1\cap L_2$). Luego ambos est\'an en $\Pclass$.
\end{proof}

\begin{porque}
La clave es que ``correr dos algoritmos polin\'omicos uno detr\'as de otro'' sigue
siendo polin\'omico: la \emph{suma} de dos polinomios es un polinomio. Por eso
$\Pclass$ aguanta combinaciones finitas sin salirse de la clase.
\end{porque}

\begin{concepto}[Cierre bajo complemento de $\Pclass$ --- se usar\'a en el Ej. 9]
$\Pclass$ tambi\'en es cerrada para el \textbf{complemento}: si decides $L$ en tiempo
polin\'omico de forma determinista, basta \emph{intercambiar} las respuestas
``S\'i''/``No'' para decidir $\overline{L}$ igual de r\'apido. Esto funciona porque la
m\'aquina es determinista y \emph{siempre para} con una respuesta clara. Para $\NP$
\emph{no} se sabe hacer esto (de ah\'i la pregunta $\NP$ vs $\coNP$).
\end{concepto}

% ==================================================================
\section{Cierre de clases de funciones por composici\'on polin\'omica}
% ==================================================================

\begin{concepto}[De qu\'e va realmente este ejercicio]
Tenemos \emph{familias} de funciones (por ejemplo ``todos los $n^k$''). Nos
preguntamos: si compongo una funci\'on de la familia con un polinomio,
?`el resultado sigue ``cabiendo'' (en orden de magnitud) dentro de la misma familia?
Hay dos formas de componer:
\begin{itemize}[nosep]
  \item \textbf{Por la izquierda:} el polinomio va \emph{fuera}: $p(f(n))$. Como
        basta tomar $p(n)=n^l$, esto es $f(n)^{l}$ --- ``elevar $f$ a una potencia''.
  \item \textbf{Por la derecha:} el polinomio va \emph{dentro}: $f(p(n))=f(n^{l})$
        --- ``meter $n^l$ como argumento de $f$''.
\end{itemize}
\textbf{Cerrada} = el resultado est\'a acotado ($=O(\cdot)$) por \emph{alguna}
funci\'on $g$ de la misma familia.
\end{concepto}

\begin{concepto}[Por qu\'e basta usar $p(n)=n^{l}$]
El enunciado lo aclara: como hablamos de \emph{\'ordenes} de complejidad, cualquier
polinomio $p(n)=a_k n^k+\cdots$ cumple $p(n)=O(n^k)$, es decir est\'a dominado por su
t\'ermino de mayor grado. As\'i que el ``peor'' polinomio, a efectos de orden, es el
mon\'omio $n^l$. Por eso solo necesitamos analizar $f(n)^l$ y $f(n^l)$.
\end{concepto}

A continuaci\'on analizamos las seis familias. Resumimos al final en una tabla.

% ---------------------------------------------------------------
\subsection*{(a) $\mathcal{C}=\{n^{k}: k>0\}$ (potencias de $n$)}

Sea $f(n)=n^{k}\in\mathcal{C}$.
\paragraph{Izquierda.} $p(f(n)) = (n^{k})^{l} = n^{kl}$, y como $kl>0$ se tiene
$n^{kl}\in\mathcal{C}$. \textbf{Cerrada por la izquierda.}
\paragraph{Derecha.} $f(p(n)) = (n^{l})^{k} = n^{lk}\in\mathcal{C}$.
\textbf{Cerrada por la derecha.}
\medskip\noindent\fbox{(a) cerrada por ambos lados.}

\begin{porque}
Funciona porque la familia est\'a parametrizada por el \emph{exponente}, y al
elevar/sustituir solo \textbf{multiplicamos exponentes} ($k\cdot l$), que sigue siendo
un exponente positivo v\'alido. La familia es ``inmune'' a estas operaciones.
\end{porque}

% ---------------------------------------------------------------
\subsection*{(b) $\mathcal{C}=\{n\cdot k: k>0\}$ (funciones lineales)}

\begin{cuidado}
Aqu\'i $k$ es un \textbf{factor} (la pendiente), no un exponente: $f(n)=k\,n$ es una
recta. No confundir con el caso (a). Todas las funciones de esta familia son
\emph{lineales}.
\end{cuidado}

\paragraph{Izquierda.} $p(f(n)) = (k\,n)^{l} = k^{l}\,n^{l}$. Para $l\geq 2$ esto es
$\Theta(n^{l})$, un crecimiento \emph{superlineal} que \textbf{no} puede acotarse por
ninguna funci\'on lineal $g(n)=k'n$ (pues $n^{l}/n=n^{l-1}\to\infty$). \textbf{No
cerrada por la izquierda.}
\paragraph{Derecha.} $f(p(n)) = k\,n^{l}$. Igual: para $l\geq 2$ es $\Theta(n^{l})$,
no acotable por una recta. \textbf{No cerrada por la derecha.}
\medskip\noindent\fbox{(b) no cerrada por ning\'un lado.}

\begin{porque}
Las rectas son ``demasiado r\'igidas'': solo tienen libertad en la pendiente, no en el
exponente, que est\'a fijo en $1$. En cuanto el polinomio introduce un exponente
$l\ge 2$, salimos de las rectas y caemos en cuadr\'aticas/c\'ubicas, que ninguna recta
puede dominar.
\end{porque}

% ---------------------------------------------------------------
\subsection*{(c) $\mathcal{C}=\{k^{n}: k>0\}$ (exponenciales de base constante)}

Sea $f(n)=k^{n}$ (interesante para $k\ge 2$).
\paragraph{Izquierda.} $p(f(n)) = (k^{n})^{l} = k^{ln} = (k^{l})^{n}$. Con
$k'=k^{l}$ obtenemos $(k')^{n}\in\mathcal{C}$. \textbf{Cerrada por la izquierda.}
\paragraph{Derecha.} $f(p(n)) = k^{\,n^{l}}$. Para $l\geq 2$ crece como $k^{n^2}$,
much\'isimo m\'as r\'apido que cualquier $(k')^{n}$:
\[
\frac{k^{\,n^{l}}}{(k')^{\,n}} = k^{\,n^{l}-n\log_k k'}\longrightarrow\infty .
\]
\textbf{No cerrada por la derecha.}
\medskip\noindent\fbox{(c) cerrada por la izquierda, no por la derecha.}

\begin{porque}
Por la izquierda solo cambiamos la \emph{base} ($k\to k^l$): seguimos en una
exponencial ``simple''. Por la derecha cambiamos el \emph{exponente} de $n$ a $n^l$:
pasamos de crecimiento $k^n$ a $k^{n^2}$, que es una bestia much\'isimo mayor que
ninguna $k'^{\,n}$ alcanza. La familia no es lo bastante amplia para contenerla.
\end{porque}

% ---------------------------------------------------------------
\subsection*{(d) $\mathcal{C}=\{2^{\,n^{k}}: k>0\}$ (exponencial con exponente polin\'omico)}

Sea $f(n)=2^{\,n^{k}}$.
\paragraph{Izquierda.} $p(f(n)) = (2^{\,n^{k}})^{l} = 2^{\,l\,n^{k}}$. Como $l$ es
constante, para $n$ grande $l\,n^{k}\leq n^{k+1}$, luego
$2^{\,l\,n^{k}} \leq 2^{\,n^{k+1}}=O(2^{\,n^{k+1}})$ y $2^{\,n^{k+1}}\in\mathcal{C}$.
\textbf{Cerrada por la izquierda.}
\paragraph{Derecha.} $f(p(n)) = 2^{\,(n^{l})^{k}} = 2^{\,n^{lk}}\in\mathcal{C}$
(pues $lk>0$). \textbf{Cerrada por la derecha.}
\medskip\noindent\fbox{(d) cerrada por ambos lados.}

\begin{porque}
Esta familia es ``ancha'' en el exponente (cualquier $n^k$ vale ah\'i arriba), as\'i
que tanto multiplicar el exponente por una constante ($l\,n^k$, que absorbemos subiendo
a $n^{k+1}$) como elevarlo ($n^{lk}$) nos deja \emph{dentro}. Compara con (c): la
diferencia es que aqu\'i el exponente ya es \emph{polin\'omico en $n$}, no constante.
\end{porque}

% ---------------------------------------------------------------
\subsection*{(e) $\mathcal{C}=\{\log^{k}(n): k>0\}$ (potencias de logaritmo)}

Sea $f(n)=(\log n)^{k}=\log^{k} n$ (el redondeo a entero no afecta al orden).
\paragraph{Izquierda.} $p(f(n)) = (\log^{k} n)^{l} = \log^{kl} n\in\mathcal{C}$, pues
$kl>0$. \textbf{Cerrada por la izquierda.}
\paragraph{Derecha.} $f(p(n)) = \log^{k}(n^{l}) = (l\log n)^{k} = l^{k}\log^{k} n
= O(\log^{k} n)$, y $\log^{k} n\in\mathcal{C}$. \textbf{Cerrada por la derecha.}
\medskip\noindent\fbox{(e) cerrada por ambos lados.}

\begin{concepto}[La identidad clave: $\log(n^{l}) = l\log n$]
Una propiedad b\'asica del logaritmo: el exponente sale como factor,
$\log(n^l)=l\log n$. Por eso ``meter $n^l$ dentro del log'' solo a\~nade un factor
constante $l$ (que el $O(\cdot)$ se traga), y por eso (e) y (f) son cerradas \emph{por
la derecha}. Es el mismo mecanismo que hace los logaritmos tan robustos.
\end{concepto}

% ---------------------------------------------------------------
\subsection*{(f) $\mathcal{C}=\{\log n\}$ (un \'unico elemento)}

\begin{cuidado}
Esta familia tiene \textbf{una sola} funci\'on, $f(n)=\log n$. Por tanto la \'unica
$g$ candidata para dominar el resultado es la propia $\log n$. Esto la hace mucho
m\'as fr\'agil que (e): no podemos ``subir el \'indice $k$'' porque no hay m\'as
elementos.
\end{cuidado}

\paragraph{Izquierda.} $p(f(n)) = (\log n)^{l} = \log^{l} n$. Para $l\geq 2$,
$\log^{l} n / \log n = \log^{l-1} n\to\infty$, as\'i que $\log^{l} n\neq O(\log n)$:
la \'unica $g$ disponible no lo acota. \textbf{No cerrada por la izquierda.}
\paragraph{Derecha.} $f(p(n)) = \log(n^{l}) = l\log n = O(\log n)$, dominado por la
\'unica funci\'on de la clase. \textbf{Cerrada por la derecha.}
\medskip\noindent\fbox{(f) cerrada por la derecha, no por la izquierda.}

\begin{porque}
Compara (e) y (f): por la izquierda obtenemos $\log^l n$ en ambos casos, pero en (e)
\emph{esa} funci\'on pertenece a la familia (hay un $\log^k$ para cada $k$), mientras
que en (f) no existe nadie m\'as que $\log n$ para acotarla. La moraleja: el cierre
no depende solo del resultado de la operaci\'on, sino de \emph{qu\'e funciones tiene
la familia disponibles} para dominarlo.
\end{porque}

% ---------------------------------------------------------------
\subsection*{Resumen}

\begin{center}
\renewcommand{\arraystretch}{1.35}
\begin{tabular}{|l|c|c|}
\hline
\textbf{Clase} & \textbf{Izquierda} $f(n)^{l}$ & \textbf{Derecha} $f(n^{l})$\\
\hline
(a) $\{n^{k}\}$        & S\'i & S\'i \\
(b) $\{k\,n\}$         & No   & No   \\
(c) $\{k^{n}\}$        & S\'i & No   \\
(d) $\{2^{\,n^{k}}\}$  & S\'i & S\'i \\
(e) $\{\log^{k} n\}$   & S\'i & S\'i \\
(f) $\{\log n\}$       & No   & S\'i \\
\hline
\end{tabular}
\end{center}

\begin{nota}
Patr\'on conceptual: la composici\'on \emph{por la izquierda} eleva $f$ a una
potencia; sobrevive cuando la familia es estable al elevar a potencias constantes
(familias param\'etricas en el exponente o el \'indice del logaritmo). La composici\'on
\emph{por la derecha} mete $n^l$ dentro de $f$; sobrevive cuando ``elevar el argumento
a una potencia'' equivale a moverse dentro de la familia. Las familias con un solo
grado de libertad mal situado (lineales, o el logaritmo aislado) fallan en el lado que
ataca ese grado de libertad.
\end{nota}

% ==================================================================
\section{$L=\{wcw : w\in\{0,1\}^{*}\}$ se reconoce en espacio $\log n$}
% ==================================================================

\begin{concepto}[Qu\'e pide el problema]
$L=\{wcw\}$ son las cadenas formadas por una palabra $w$ de ceros y unos, una marca
$c$ en medio, y \emph{la misma} palabra $w$ otra vez. Ejemplos: $0c0$, $101c101$ son
de $L$; $10c01$ no (las mitades difieren). Queremos decidir la pertenencia usando
solo $O(\log n)$ memoria de trabajo.
\end{concepto}

\begin{cuidado}
La tentaci\'on es \emph{copiar} la primera mitad $w$ a la cinta de trabajo y luego
compararla. Pero copiar $w$ cuesta $|w|\approx n/2$ espacio: \textbf{lineal}, no
logar\'itmico. Hay que evitarlo a toda costa.
\end{cuidado}

\begin{teorema}
$L=\{wcw : w\in\{0,1\}^{*}\}$ pertenece a $\Lclass=\SPACE(\log n)$.
\end{teorema}

\begin{idea}
En vez de copiar, comparamos las dos mitades \emph{posici\'on a posici\'on} usando
\textbf{punteros} (\'indices). Comparar el car\'acter $i$ de la izquierda con el
car\'acter $i$ de la derecha solo requiere recordar el n\'umero $i$, que cabe en
$O(\log n)$ bits. Aqu\'i es esencial que la cinta de entrada sea de solo lectura y
podamos movernos por ella libremente (concepto base).
\end{idea}

\begin{proof}
Modelo: cinta de entrada de solo lectura + cinta de trabajo (la que se mide). Entrada
de longitud $n$.

\textbf{Fase 1 (formato).} Recorremos la entrada contando cu\'antas $c$ hay y
anotando la posici\'on $m$ de la (\'unica) $c$. El contador y $m$ son $\le n$, luego
$O(\log n)$ bits. Si el n\'umero de $c$ no es exactamente $1$, rechazar. Tras esto,
la izquierda es $u=x_1\cdots x_{m-1}$ y la derecha $v=x_{m+1}\cdots x_n$.

\textbf{Fase 2 (longitudes).} Para que sea $wcw$ debe cumplirse $|u|=|v|$, o sea
$m-1=n-m$. Se comprueba con aritm\'etica sobre $m$ y $n$, en $O(\log n)$. Si no,
rechazar.

\textbf{Fase 3 (comparaci\'on).} Para $i=1,2,\dots,m-1$:
\begin{itemize}[nosep]
  \item mover la cabeza de lectura a la posici\'on $i$ y leer $x_i$ (un bit);
  \item mover la cabeza a la posici\'on $m+i$ y leer $x_{m+i}$;
  \item si $x_i\neq x_{m+i}$, rechazar.
\end{itemize}
Si todas coinciden, aceptar.

\textbf{Espacio.} En la cinta de trabajo solo hay, en cada momento, un contador
(el \'indice $i$, o $m$, o el conteo de $c$) y uno o dos bits le\'idos. Todo \'indice
es $\le n$, luego $O(\log n)$ bits; el n\'umero de variables es constante. Total:
$O(\log n)$.
\end{proof}

\begin{porque}
?`Por qu\'e ``$O(\log n)$'' y no contar el tiempo? El tiempo es $O(n^2)$ (por los
reposicionamientos de la cabeza), pero el ejercicio pide \emph{espacio}, y en espacio
ganamos much\'isimo: pasamos de memorizar media palabra ($n/2$ bits) a memorizar un
\'indice ($\log n$ bits). El truco ``releer en lugar de almacenar'' reaparecer\'a en
los ejercicios 10--12.
\end{porque}

% ==================================================================
\section{Si $L\in\Pclass$ entonces $\Lstar\in\Pclass$}
% ==================================================================

\begin{concepto}[La estrella de Kleene $\Lstar$]
$\Lstar$ es el conjunto de cadenas que se obtienen \emph{concatenando cero o m\'as}
trozos de $L$. Es decir, $x\in\Lstar$ si $x$ se puede partir como
$x=w_1 w_2\cdots w_k$ con cada $w_j\in L$ (y la cadena vac\'ia $\varepsilon$ siempre
est\'a, con $k=0$). Ejemplo: si $L=\{ab,c\}$, entonces $abcab\in\Lstar$ (partido como
$ab\cdot c\cdot ab$).
\end{concepto}

\begin{teorema}
Si $L\in\Pclass$, entonces $\Lstar\in\Pclass$.
\end{teorema}

\begin{idea}
El problema es ``?`existe \emph{alguna} forma de cortar $x$ en trozos que est\'en
todos en $L$?''. Probar todas las formas de cortar ser\'ia exponencial. Usamos
\textbf{programaci\'on din\'amica}: resolvemos el problema para los prefijos cada vez
m\'as largos, reutilizando lo ya calculado.
\end{idea}

\begin{concepto}[Programaci\'on din\'amica, en una frase]
T\'ecnica que rompe un problema en subproblemas que se \emph{solapan}, los resuelve
una sola vez y guarda las respuestas en una tabla para no repetir trabajo. Aqu\'i el
subproblema es: ``?`el prefijo $x[1..j]$ pertenece a $\Lstar$?''.
\end{concepto}

\begin{proof}
Sea $M$ una m\'aquina que decide $L$ en tiempo $O(n^{c})$. Escribimos
$x[i..j]=x_i\cdots x_j$ (subcadena) y por convenio $x[i..i-1]=\varepsilon$.

Definimos el vector booleano $A[0..n]$ con
\[
A[j] = \text{Verdadero} \iff x[1..j]\in\Lstar .
\]
Recurrencia (mirando d\'onde empieza el \emph{\'ultimo} trozo):
\[
A[0]=\text{Verdadero}, \qquad
A[j]=\bigvee_{0\le i<j}\Big(A[i]\ \wedge\ \big(x[i+1..j]\in L\big)\Big).
\]

\textbf{Algoritmo.} Calcular $A[0],A[1],\dots,A[n]$ en orden creciente; para cada par
$(i,j)$ se invoca $M$ sobre $x[i+1..j]$. Aceptar si $A[n]$ es Verdadero.

\textbf{Correcci\'on.} Por inducci\'on sobre $j$: $A[j]$ es Verdadero exactamente
cuando $x[1..j]$ se puede descomponer en factores de $L$, que es la definici\'on de
pertenecer a $\Lstar$. El caso base $A[0]$ es $\varepsilon\in\Lstar$.

\textbf{Complejidad.} Hay $O(n^2)$ pares $(i,j)$; cada uno cuesta una llamada a $M$,
$O(n^c)$. Total $O(n^{c+2})$, polin\'omico. Luego $\Lstar\in\Pclass$.
\end{proof}

\begin{porque}
?`Por qu\'e mirar ``d\'onde empieza el \'ultimo trozo''? Porque si conocemos ese punto
de corte $i$, el problema se parte limpiamente en dos piezas \emph{independientes}:
``?`$x[1..i]$ ya estaba bien partido?'' (eso es $A[i]$, ya calculado) y ``?`el resto
$x[i+1..j]$ es un \'unico trozo de $L$?'' (una llamada a $M$). El ``o'' ($\bigvee$)
recorre todos los posibles puntos de corte: con que \emph{uno} funcione, basta.
\end{porque}

% ==================================================================
\section{Si $L\in\NP$ entonces $\Lstar\in\NP$}
% ==================================================================

\begin{teorema}
Si $L\in\NP$, entonces $\Lstar\in\NP$.
\end{teorema}

\begin{idea}
Usamos la caracterizaci\'on por \emph{certificado} (concepto base). Para convencer a
alguien de que $x\in\Lstar$ basta darle: (1) d\'onde est\'an los cortes que parten
$x$ en trozos, y (2) para cada trozo, el certificado de que ese trozo est\'a en $L$.
Todo eso es corto y se comprueba r\'apido.
\end{idea}

\begin{proof}[Demostraci\'on (versi\'on verificador)]
Como $L\in\NP$, hay un verificador $V$ y un polinomio $q$ con
$w\in L\iff \exists c_w\,(|c_w|\le q(|w|)):\ V(w,c_w)=1$, y $V$ corre en tiempo
polin\'omico.

Definimos un verificador $V^{*}$ para $\Lstar$. Dado $x$ ($|x|=n$), el certificado
propuesto es
\[
\big(\, 0=i_0 < i_1 < \cdots < i_k = n,\ \ c_1,\dots,c_k \,\big),
\]
donde los $i_j$ marcan los cortes (el trozo $j$ es $w_j=x[i_{j-1}+1..i_j]$) y $c_j$
certifica $w_j\in L$. Si $x=\varepsilon$ se acepta ($k=0$). $V^{*}$:
\begin{enumerate}[nosep]
  \item comprueba que $0=i_0<i_1<\cdots<i_k=n$ (cortes v\'alidos que cubren todo $x$);
  \item para cada $j$, extrae $w_j$ y verifica $V(w_j,c_j)=1$;
  \item acepta si todo va bien.
\end{enumerate}

\textbf{Correcci\'on.} $x\in\Lstar$ syss existe una partici\'on en trozos de $L$, lo
cual equivale a la existencia de cortes v\'alidos y certificados aceptados.

\textbf{Tama\~no del certificado.} A lo sumo $k\le n$ cortes, cada uno un \'indice
$\le n$ ($O(n\log n)$ bits); cada $c_j$ mide $\le q(n)$ y, como $\sum_j|w_j|=n$, la
suma total de certificados es $\le n\cdot q(n)$: polin\'omica.

\textbf{Tiempo.} Comprobar cortes es lineal; ejecutar $V$ en cada trozo es
polin\'omico y se hace $\le n$ veces. Polin\'omico en total. Luego $\Lstar\in\NP$.
\end{proof}

\begin{concepto}[La diferencia entre el Ej.\ 6 y el Ej.\ 7]
Son el ``mismo'' problema ($\Lstar$) pero con distinta hip\'otesis sobre $L$:
\begin{itemize}[nosep]
  \item Ej.\ 6 ($L\in\Pclass$): podemos \emph{comprobar} ``$x[i+1..j]\in L$''
        directamente, as\'i que basta la programaci\'on din\'amica determinista.
  \item Ej.\ 7 ($L\in\NP$): ya no podemos comprobarlo de golpe, pero \emph{s\'i}
        verificar un certificado. As\'i que el certificado global de $\Lstar$ = los
        cortes $+$ un certificado por trozo.
\end{itemize}
De hecho se podr\'ia hacer el 7 con la misma programaci\'on din\'amica del 6,
sustituyendo cada test ``$\in L$'' por ``adivinar y verificar el certificado del
trozo''. Ambos enfoques son v\'alidos.
\end{concepto}

% ==================================================================
\section{$\NP$ es cerrada para uni\'on e intersecci\'on}
% ==================================================================

\begin{teorema}
Si $L_1,L_2\in\NP$, entonces $L_1\cup L_2\in\NP$ y $L_1\cap L_2\in\NP$.
\end{teorema}

\begin{idea}
Como en el Ej.\ 3, pero ahora con certificados en lugar de m\'aquinas deterministas.
La pregunta clave en cada caso es: ``?`qu\'e pista (certificado) basta para convencer
de que $x$ est\'a en la uni\'on / intersecci\'on?''.
\end{idea}

\begin{proof}
Sean $V_1,V_2$ verificadores polin\'omicos y $q_1,q_2$ polinomios con
$x\in L_i\iff \exists c_i\,(|c_i|\le q_i(|x|)):\ V_i(x,c_i)=1$.

\paragraph{Uni\'on.} El certificado es un par (etiqueta, testigo): un bit
$b\in\{1,2\}$ que dice ``a cu\'al de los dos pertenece $x$'' m\'as el certificado
$c_b$ correspondiente. El verificador $V_\cup(x,(b,c))$ ejecuta $V_b(x,c)$. Entonces
\[
x\in L_1\cup L_2 \iff (\exists c_1: V_1(x,c_1)=1)\ \text{o}\ (\exists c_2:
V_2(x,c_2)=1) \iff \exists (b,c): V_\cup(x,(b,c))=1 .
\]

\paragraph{Intersecci\'on.} El certificado es el \emph{par} $(c_1,c_2)$ de \emph{ambos}
testigos. $V_\cap(x,(c_1,c_2))$ ejecuta $V_1(x,c_1)$ \textbf{y} $V_2(x,c_2)$ y acepta
si ambos aceptan. Entonces
\[
x\in L_1\cap L_2 \iff (\exists c_1: V_1(x,c_1)=1)\ \text{y}\ (\exists c_2:
V_2(x,c_2)=1)\iff \exists(c_1,c_2): V_\cap(x,(c_1,c_2))=1 .
\]
En ambos casos el certificado es de tama\~no polin\'omico y el verificador corre en
tiempo polin\'omico. Luego $L_1\cup L_2,\ L_1\cap L_2\in\NP$.
\end{proof}

\begin{cuidado}
En la intersecci\'on, el paso ``$\exists c_1\,\exists c_2 \equiv \exists (c_1,c_2)$''
es l\'icito porque los dos testigos son \textbf{independientes}: elegir $c_1$ no
condiciona la b\'usqueda de $c_2$. Esto NO funcionar\'ia si la operaci\'on requiriese
cuantificadores alternados ($\exists\dots\forall\dots$); ah\'i es justo donde $\NP$
deja de ser suficiente y aparecen clases superiores. Pero uni\'on e intersecci\'on
solo usan ``$\exists$'', as\'i que $\NP$ aguanta.
\end{cuidado}

% ==================================================================
\section{Si $\NP\neq\coNP$ entonces $\Pclass\neq\NP$}
% ==================================================================

\begin{concepto}[Demostrar por el contrarrec\'iproco]
Probar ``$A\Rightarrow B$'' es \emph{exactamente} lo mismo que probar
``$\text{no }B\Rightarrow \text{no }A$'' (el \textbf{contrarrec\'iproco}). A veces una
direcci\'on es mucho m\'as f\'acil. Aqu\'i, en vez de ``$\NP\neq\coNP \Rightarrow
\Pclass\neq\NP$'', probamos lo equivalente y m\'as manejable:
``$\Pclass=\NP \Rightarrow \NP=\coNP$''.
\end{concepto}

\begin{teorema}
Si $\NP\neq\coNP$, entonces $\Pclass\neq\NP$.
\end{teorema}

\begin{proof}
Probamos el contrarrec\'iproco: $\Pclass=\NP \Rightarrow \NP=\coNP$.

Supongamos $\Pclass=\NP$. Recordemos (concepto del Ej.\ 3) que $\Pclass$ es cerrada
para complemento: si $A\in\Pclass$, intercambiando aceptaci\'on/rechazo decidimos
$\overline{A}$ igual de r\'apido, luego $\overline{A}\in\Pclass$.

Veamos la doble inclusi\'on $\NP=\coNP$.

\textbf{$\NP\subseteq\coNP$:} Sea $A\in\NP$. Por hip\'otesis $A\in\Pclass$, y como
$\Pclass$ es cerrada para complemento, $\overline{A}\in\Pclass=\NP$. Pero
``$\overline{A}\in\NP$'' es justo la definici\'on de ``$A\in\coNP$''. Luego
$A\in\coNP$.

\textbf{$\coNP\subseteq\NP$:} Sea $A\in\coNP$, es decir $\overline{A}\in\NP$. Por
hip\'otesis $\overline{A}\in\Pclass$, y por cierre bajo complemento
$A\in\Pclass=\NP$. Luego $A\in\NP$.

De ambas, $\NP=\coNP$. Tomando contrarrec\'iproco:
$\NP\neq\coNP \Rightarrow \Pclass\neq\NP$.
\end{proof}

\begin{porque}
El \emph{coraz\'on} del argumento es la asimetr\'ia: $\Pclass$ \textbf{s\'i} es cerrada
para complemento (porque sus m\'aquinas son deterministas y siempre paran), pero de
$\NP$ \emph{no} se sabe. Si $\Pclass$ y $\NP$ fueran iguales, $\NP$ heredar\'ia esa
propiedad de $\Pclass$ y colapsar\'ia con su complementaria $\coNP$. Que la mayor\'ia
de expertos crea $\NP\neq\coNP$ es, por esta v\'ia, otra raz\'on para creer
$\Pclass\neq\NP$.
\end{porque}

\begin{concepto}[Definici\'on de $\coNP$, recordatorio aplicado]
$A\in\coNP$ significa, por definici\'on, que su complementario $\overline{A}$ est\'a
en $\NP$. Por eso en la prueba ``traducimos'' constantemente entre ``$A\in\coNP$'' y
``$\overline{A}\in\NP$'': son la misma afirmaci\'on. Tenerlo claro es lo que hace la
prueba casi mec\'anica.
\end{concepto}

% ==================================================================
\section{Par\'entesis bien emparejados est\'a en $\Lclass$}
% ==================================================================

\begin{teorema}
El lenguaje de cadenas de par\'entesis $\{(,)\}$ correctamente emparejados y anidados
pertenece a $\Lclass=\SPACE(\log n)$.
\end{teorema}

\begin{idea}
Con \textbf{un solo} tipo de par\'entesis no hace falta una pila: basta un
\textbf{contador} de profundidad (cu\'antos par\'entesis hay abiertos en este
momento). La regla de buena formaci\'on es local: ``el contador nunca baja de cero y
acaba en cero''.
\end{idea}

\begin{proof}
Entrada $x=x_1\cdots x_n$ sobre $\{(,)\}$. Mantenemos un \'unico contador $b$
(``balance''), $b\leftarrow 0$. Recorremos de izquierda a derecha:
\begin{itemize}[nosep]
  \item si $x_i = $ ``\texttt{(}'', hacer $b\leftarrow b+1$;
  \item si $x_i = $ ``\texttt{)}'', hacer $b\leftarrow b-1$;
  \item si en alg\'un momento $b<0$, \textbf{rechazar}.
\end{itemize}
Al terminar, \textbf{aceptar} syss $b=0$.

\textbf{Correcci\'on.} Una cadena de un solo tipo est\'a bien formada syss en todo
prefijo hay tantos o m\'as ``\texttt{(}'' que ``\texttt{)}'' (nunca se cierra de
m\'as) y al final hay igual n\'umero (todo lo abierto se cierra). El contador mide
exactamente \#\texttt{(} $-$ \#\texttt{)} del prefijo le\'ido: ``$b\ge 0$ siempre'' y
``$b=0$ al final'' son justo esas dos condiciones. As\'i, \texttt{(()())} y
\texttt{((()))} se aceptan; \texttt{)()(} (balance $-1$ al primer s\'imbolo) y
\texttt{(()))(} (balance llega a $-1$) se rechazan.

\textbf{Espacio.} El contador cumple $0\le b\le n$, luego cabe en
$\lceil\log_2(n+1)\rceil=O(\log n)$ bits. No se usa nada m\'as. Est\'a en $\Lclass$.
\end{proof}

\begin{porque}
La raz\'on por la que un contador basta es que, con un solo tipo, \emph{no importa
cu\'ales} par\'entesis se emparejan, solo \emph{cu\'antos} hay abiertos. Toda la
informaci\'on de una pila de altura $h$ se resume en el n\'umero $h$. Y un n\'umero
$\le n$ cuesta solo $\log n$ bits. En cuanto haya \emph{dos} tipos (Ej.\ 11) esto se
rompe: ya importar\'a el orden.
\end{porque}

% ==================================================================
\section{Dos tipos de par\'entesis}
% ==================================================================

\begin{teorema}
El lenguaje de cadenas bien formadas con dos tipos de par\'entesis, $(,)$ y $[,]$,
tambi\'en pertenece a $\Lclass$.
\end{teorema}

\begin{concepto}[Por qu\'e dos tipos complican las cosas: la pila]
Con dos tipos importa el \emph{orden} de cierre: \texttt{([]())} es v\'alida pero
\texttt{([)]} no (se cierra un ``\texttt{)}'' cuando lo \'ultimo abierto era un
``\texttt{[}''). La estructura natural para llevar esto es una \textbf{pila}: cada
apertura se apila, cada cierre debe casar con el \'ultimo apilado (el del tope).
Pero una pila puede crecer hasta $\Theta(n)$ s\'imbolos: ser\'ia espacio
\emph{lineal}, no logar\'itmico.
\end{concepto}

\begin{idea}
El truco para mantenernos en $O(\log n)$ es \textbf{no almacenar la pila}: para cada
cierre, \emph{recalculamos} cu\'al es su apertura asociada recorriendo de nuevo la
entrada con contadores. Pagamos m\'as tiempo a cambio de casi nada de espacio.
(Aprovechamos que la cinta de entrada es de solo lectura y la podemos releer.)
\end{idea}

\begin{proof}
\textbf{Validez:} la cadena es correcta syss (1) est\'a balanceada \emph{ignorando el
tipo} (contador global nunca negativo, cero al final --- como en el Ej.\ 10), y (2)
cada cierre casa en \emph{tipo} con la apertura que lo empareja en el anidamiento.

\textbf{Emparejar sin pila.} Para un cierre en posici\'on $j$, su apertura asociada es
la posici\'on $i<j$ tal que el segmento intermedio $x[i+1..j-1]$ est\'a balanceado.
Operativamente, en $O(\log n)$ espacio: recorremos $x$ con el contador de balance $b$
(verificando de paso la condici\'on 1). Cada vez que hallamos un cierre en posici\'on
$j$, lanzamos un segundo recorrido desde el principio para localizar esa apertura
$i$ (la \'ultima posici\'on $<j$ desde la que el sub-balance hasta $j-1$ es $0$), y
comprobamos que $x_i$ y $x_j$ son del mismo tipo. Si no, rechazar.

Todo el almacenamiento son unos pocos contadores acotados por $n$ ($j$, la candidata
$i$, los balances), cada uno $O(\log n)$ bits; no guardamos la pila completa, la
recomputamos. Espacio total: $O(\log n)$. (Tiempo: $O(n^2)$.) Luego est\'a en
$\Lclass$.

\medskip
\textbf{Segunda parte: leyendo de izquierda a derecha SIN volver atr\'as, requieren
$O(n)$ de memoria.}
\end{proof}

\begin{concepto}[Modelo ``online'': una pasada, sin retroceso]
Ahora cambiamos las reglas: la entrada se lee \textbf{una sola vez}, de izquierda a
derecha, sin poder releer (como un flujo que pasa y no vuelve). En este modelo ya
\emph{no} podemos ``recomputar la pila'' como antes (eso requer\'ia releer). La
pregunta es cu\'anta memoria necesitamos \emph{obligatoriamente}.
\end{concepto}

\begin{concepto}[Argumento de distinguibilidad (``fooling set'')]
T\'ecnica para demostrar \emph{cotas inferiores} de memoria. Idea: si encontramos
muchos prefijos distintos que la m\'aquina \textbf{tiene} que tratar de forma
diferente (porque su continuaci\'on correcta difiere), entonces tras leerlos debe
estar en \emph{configuraciones de memoria distintas}. Si hay $N$ prefijos
distinguibles, hacen falta $\ge \log_2 N$ bits de memoria. Es el mismo principio que
dice ``para distinguir $N$ cosas necesitas $\log_2 N$ bits''.
\end{concepto}

\begin{proof}[Demostraci\'on de la cota $\Omega(n)$]
Consideremos los $2^{m}$ prefijos de $m$ aperturas, cada una de tipo redondo o
cuadrado: $u_S=s_1\cdots s_m$ con $s_t\in\{\,\texttt{(}\,,\,\texttt{[}\,\}$. Para dos
prefijos distintos $u_S\neq u_{S'}$, sea $t$ la primera posici\'on donde difieren los
tipos. La continuaci\'on que cierra \emph{en el orden correcto para $u_S$} ---el
``espejo'' $\bar s_m\cdots\bar s_1$, donde $\bar{\texttt{(}}=\texttt{)}$ y
$\bar{\texttt{[}}=\texttt{]}$--- completa $u_S$ a una cadena \textbf{v\'alida}, pero
aplicada a $u_{S'}$ produce un \textbf{desajuste de tipo}, dando cadena
\textbf{inv\'alida}.

Por tanto la m\'aquina debe quedar en estados de memoria \emph{distintos} tras leer
$u_S$ y $u_{S'}$: si coincidieran, al recibir la misma continuaci\'on se comportar\'ia
igual con ambas (sin poder releer el prefijo), aceptando o rechazando las dos a la
vez, contradicci\'on.

As\'i hay $\ge 2^{m}$ configuraciones de memoria distinguibles tras leer $m$
s\'imbolos, lo que exige $\ge\log_2(2^m)=m$ bits de memoria. Como $m$ puede ser
$\Theta(n)$, la memoria necesaria es $\Omega(n)$ (lineal).
\end{proof}

\begin{porque}
?`Por qu\'e usamos el ``espejo'' como continuaci\'on? Porque es la \'unica forma de
cerrar bien un anidamiento: lo \'ultimo que se abre es lo primero que se cierra. Al
fijar esa continuaci\'on para $u_S$, forzamos que cualquier $u_{S'}$ con tipos
distintos falle. As\'i fabricamos a prop\'osito un mont\'on de pares que la m\'aquina no
puede confundir, y de ah\'i sale la cota.
\end{porque}

\begin{cuidado}
\textbf{No te confundas:} la primera parte (en $\Lclass$, $O(\log n)$) y la segunda
(requiere $O(n)$) \emph{no} se contradicen. Hablan de \textbf{modelos distintos}: la
primera \emph{permite releer} la entrada (acceso de solo lectura con
reposicionamiento), la segunda \emph{no}. La cota logar\'itmica se consigue
precisamente \emph{gracias} a poder releer y recomputar la pila; si lo prohibimos, hay
que memorizar la secuencia de tipos, que son $\Theta(n)$ bits.
\end{cuidado}

% ==================================================================
\section{El pal\'indromo requiere $O(n)$ sin retroceder}
% ==================================================================

\begin{concepto}[Pal\'indromo]
Una cadena $w$ es \textbf{pal\'indromo} si se lee igual al derecho y al rev\'es:
$w=w^{R}$ (donde $w^R$ es $w$ invertida). Ejemplos: \texttt{0110}, \texttt{10101}. El
problema es decidir si la entrada es un pal\'indromo, en el modelo de una sola pasada
sin retroceso.
\end{concepto}

\begin{teorema}
Reconocer pal\'indromos requiere memoria $\Omega(n)$ (lineal) si la entrada se lee de
izquierda a derecha sin poder volver atr\'as.
\end{teorema}

\begin{idea}
Mismo \emph{argumento de distinguibilidad} del Ej.\ 11. Exhibimos exponencialmente
muchos prefijos que la m\'aquina debe distinguir (su continuaci\'on correcta difiere),
forzando $2^m$ configuraciones de memoria y por tanto $\Omega(m)=\Omega(n)$ bits.
\end{idea}

\begin{proof}
Sea $L_{\text{pal}}=\{w:w=w^R\}$ sobre $\{0,1\}$, y $M$ una m\'aquina que lee la
entrada una vez, de izquierda a derecha, con memoria de trabajo $s(n)$.

Fijado $m$, tomamos los $2^m$ prefijos $u\in\{0,1\}^m$. Afirmamos que $M$ debe llegar
a configuraciones de memoria distintas tras leer dos prefijos distintos $u\neq u'$.

En efecto, supongamos lo contrario: $M$ alcanza la \emph{misma} configuraci\'on tras
leer $u$ y tras $u'$, con $u\neq u'$ (difieren en alguna posici\'on $t$). Usamos la
continuaci\'on $z=u^R$ (el reverso de $u$). Entonces:
\begin{itemize}[nosep]
  \item $u\,z=u\,u^R$ es pal\'indromo (forma $w w^R$): $M$ debe \textbf{aceptar}.
  \item $u'\,z=u'\,u^R$ \textbf{no} es pal\'indromo: en la posici\'on $t$ (primera
        mitad) vale $u'_t$, mientras que su sim\'etrica (segunda mitad) vale
        $u_t\neq u'_t$: $M$ debe \textbf{rechazar}.
\end{itemize}
Pero las configuraciones tras $u$ y $u'$ coinciden por hip\'otesis, y como $M$ no
puede volver atr\'as, su comportamiento al leer $z$ es id\'entico en ambos casos:
aceptar\'ia las dos o rechazar\'ia las dos. Contradicci\'on.

Luego las $2^m$ configuraciones son distintas. El n\'umero de configuraciones con $s$
celdas (control de $Q$ estados, cabeza en $s$ posiciones, alfabeto de trabajo de
tama\~no $\gamma$) es $\le Q\cdot s\cdot \gamma^{s}=2^{O(s)}$. Necesitamos
$2^{O(s)}\ge 2^m$, de donde $s=\Omega(m)$. Tomando $m=\lfloor n/2\rfloor$,
$s(n)=\Omega(n)$.
\end{proof}

\begin{concepto}[Por qu\'e ``$\#$configuraciones $=2^{O(s)}$'']
Una ``configuraci\'on de memoria'' es una foto del estado interno: qu\'e estado de
control, d\'onde est\'a la cabeza, y qu\'e contiene la cinta de trabajo. Con $s$
celdas y un alfabeto de $\gamma$ s\'imbolos, los contenidos posibles son $\gamma^s$;
multiplicado por $Q$ estados y $s$ posiciones de cabeza da $Q\,s\,\gamma^s$, que es
$2^{O(s)}$. La moraleja: \textbf{con $s$ bits de memoria solo puedes ``recordar''
$2^{O(s)}$ situaciones distintas}. Si necesitas distinguir $2^m$, te hacen falta
$s\ge\Omega(m)$ bits.
\end{concepto}

\begin{cuidado}
Compara con el Ej.\ 5: $\{wcw\}$ \emph{s\'i} est\'a en $\Lclass$, y el pal\'indromo
$\{ww^R\}$ tambi\'en lo estar\'ia \emph{con relectura} (comparando $x_i$ con
$x_{n+1-i}$ por punteros). La cota $\Omega(n)$ de este ejercicio depende
\textbf{cr\'iticamente} de ``sin volver atr\'as'': es lo que proh\'ibe los punteros y
obliga a memorizar media cadena entera. Si en un examen ves ``espacio $\log n$'' y
``$\Omega(n)$'' para problemas parecidos, casi siempre la diferencia es el modelo de
acceso a la entrada.
\end{cuidado}

% ==================================================================
\section{$\NP\neq\SPACE(n)$}
% ==================================================================

\begin{cuidado}
Este ejercicio es \emph{el m\'as dif\'icil} de la relaci\'on. Un punto clave para no
perderse: \textbf{no} vamos a decir cu\'al clase contiene a cu\'al (de hecho se
desconoce si $\NP\subseteq\SPACE(n)$). Solo demostraremos que \emph{no pueden ser la
misma}. La t\'actica: encontrar una \textbf{propiedad} que una clase tiene y la otra
no. Dos clases iguales tendr\'ian las mismas propiedades; si difieren en una, son
distintas.
\end{cuidado}

\begin{idea}
La propiedad elegida (sugerida por el enunciado) es ``ser cerrada para reducciones en
espacio logar\'itmico''. Mostraremos: \textbf{(Paso 1)} $\NP$ s\'i lo es;
\textbf{(Paso 2)} $\SPACE(n)$ no lo es. Conclusi\'on inmediata: $\NP\neq\SPACE(n)$.
\end{idea}

\begin{definicion}[cierre para reducciones en $\Lclass$]
Una clase $\mathcal{D}$ es \emph{cerrada para reducciones en espacio logar\'itmico}
si: siempre que $P_1\reduce P_2$ (reducci\'on calculable en espacio $\log$) y
$P_2\in\mathcal{D}$, entonces tambi\'en $P_1\in\mathcal{D}$.
\end{definicion}

\begin{concepto}[Qu\'e dice esta propiedad, en cristiano]
``$\mathcal{D}$ es cerrada para $\log$-reducciones'' significa: \emph{si un problema
se reduce baratamente a otro que est\'a en $\mathcal{D}$, el primero tambi\'en
est\'a en $\mathcal{D}$}. Es una forma de decir que $\mathcal{D}$ ``respeta'' la
relaci\'on de dificultad: nada que sea ``m\'as f\'acil o igual'' que un miembro de
$\mathcal{D}$ se queda fuera.
\end{concepto}

\subsection*{Paso 1: $\NP$ es cerrada para reducciones en $\Lclass$}

\begin{lema}
Si $P_1\reduce P_2$ y $P_2\in\NP$, entonces $P_1\in\NP$.
\end{lema}

\begin{proof}
Sea $f$ la funci\'on de reducci\'on, calculable en espacio $O(\log n)$ y por tanto en
tiempo polin\'omico (ver concepto siguiente); en particular $|f(x)|$ es polin\'omico.
Cumple $x\in P_1\iff f(x)\in P_2$. Como $P_2\in\NP$, sea $V_2$ su verificador
polin\'omico.

Verificador para $P_1$: dado $x$ y un certificado $c$, (1) calcular $y=f(x)$,
(2) aceptar syss $V_2(y,c)=1$. Entonces
\[
x\in P_1 \iff f(x)\in P_2 \iff \exists c\,(|c|\le q(|f(x)|)):\ V_2(f(x),c)=1,
\]
y como $|f(x)|$ es polin\'omico en $|x|$, el certificado es corto y la verificaci\'on
polin\'omica. Luego $P_1\in\NP$.
\end{proof}

\begin{concepto}[Por qu\'e ``espacio $\log$'' implica ``tiempo polin\'omico'']
Una m\'aquina que usa $O(\log n)$ espacio de trabajo tiene un n\'umero \emph{limitado}
de configuraciones distintas: a lo sumo $2^{O(\log n)}=n^{O(1)}$ (polin\'omicas). Si
una m\'aquina que siempre para no repitiera configuraci\'on (si lo hiciera, entrar\'ia
en bucle infinito), entonces da a lo sumo $n^{O(1)}$ pasos. Conclusi\'on:
$\Lclass\subseteq\Pclass$ y, en general, todo c\'omputo $\log$-espacial corre en
tiempo polin\'omico. Por eso la reducci\'on, aunque definida por espacio, produce una
salida $f(x)$ de tama\~no solo polin\'omico.
\end{concepto}

\subsection*{Paso 2: $\SPACE(n)$ \emph{no} es cerrada para reducciones en $\Lclass$}

\begin{concepto}[Diagonalizaci\'on y el Teorema de Jerarqu\'ia de Espacio]
\textbf{Diagonalizaci\'on} es la t\'ecnica (la de Cantor, la del problema de la
parada) que construye un objeto que difiere de cada elemento de una lista en al menos
un punto, garantizando que ``no est\'a en la lista''. Aplicada a complejidad, permite
construir un problema que ninguna m\'aquina con \emph{pocos} recursos puede resolver,
pero s\'i una con \emph{algo m\'as} de recursos. El resultado:
\end{concepto}

\begin{teorema}[Jerarqu\'ia de espacio; Stearns--Hartmanis--Lewis]
Si $s_1(n)=o(s_2(n))$ (y $s_2$ es ``construible en espacio'', $s_2(n)\ge\log n$),
entonces la inclusi\'on $\SPACE(s_1(n))\subsetneq\SPACE(s_2(n))$ es
\textbf{estricta}. En particular, como $n=o(n^2)$:
\[
\SPACE(n)\subsetneq\SPACE(n^2).
\]
\end{teorema}

\begin{porque}
Necesitamos este teorema porque nos da, ``gratis'', un problema concreto que est\'a en
$\SPACE(n^2)$ pero \emph{no} en $\SPACE(n)$: justo la pieza que separa una clase de la
otra. Sin \'el no tendr\'iamos de d\'onde sacar ese testigo. La demostraci\'on del
teorema es por diagonalizaci\'on (se fabrica una m\'aquina que, en espacio $n^2$,
contradice a toda m\'aquina de espacio lineal), pero para este ejercicio lo usamos
como caja negra.
\end{porque}

\begin{lema}
$\SPACE(n)$ no es cerrada para reducciones en espacio logar\'itmico.
\end{lema}

\begin{concepto}[La t\'ecnica del ``padding'' (relleno)]
\textbf{Padding} = alargar artificialmente la entrada a\~nadiendo s\'imbolos de
relleno. ?`Para qu\'e sirve? El coste en espacio se mide \emph{respecto a la longitud
de la entrada}. Si una entrada $x$ requiere espacio $|x|^2$ pero la ``estiramos'' a
longitud $N=|x|^2$, ese mismo coste $|x|^2=N$ pasa a ser \emph{lineal en $N$}. El
relleno ``diluye'' la dificultad: el problema dif\'icil se disfraza de f\'acil al
medirlo sobre una entrada mayor. Es el truco que provoca el fallo de cierre.
\end{concepto}

\begin{proof}
Por la jerarqu\'ia, $\SPACE(n)\subsetneq\SPACE(n^2)$; fijamos un lenguaje testigo
\[
A\in\SPACE(n^2)\setminus\SPACE(n)
\]
(existe por ser la inclusi\'on estricta: decidible en espacio cuadr\'atico, pero
\emph{no} en lineal).

Definimos su versi\'on rellenada. A\~nadimos un s\'imbolo nuevo $\#\notin\Sigma$:
\[
B=\{\,x\#^{\,t}\ :\ x\in A,\ t=|x|^2-|x|\,\},
\]
es decir, a cada $x\in A$ le pegamos relleno hasta que la longitud total sea
$N=|x|^2$.

\textbf{(i) $B\in\SPACE(n)$.} Dada $z$ de longitud $N$: comprobamos que tiene forma
$x\#^t$ con $t=|x|^2-|x|$ (contando, en $O(\log N)$), extraemos $x$ y ejecutamos el
decididor de $A$, que usa espacio $O(|x|^2)$. Pero $|x|^2=N$, as\'i que es $O(N)$,
\emph{lineal en la entrada $z$}. Luego $B\in\SPACE(N)=\SPACE(n)$.

\textbf{(ii) $A\reduce B$ en espacio $\log$.} La reducci\'on $f(x)=x\#^{\,|x|^2-|x|}$
se calcula en $O(\log|x|)$: copiar $x$ y escribir $|x|^2-|x|$ s\'imbolos $\#$ llevando
solo un contador. Y $x\in A\iff f(x)\in B$.

\textbf{(iii) Conclusi\'on.} Tenemos $A\reduce B$ con $B\in\SPACE(n)$. Si $\SPACE(n)$
fuera cerrada para $\log$-reducciones, deber\'ia ser $A\in\SPACE(n)$. Pero
$A\notin\SPACE(n)$ por construcci\'on. \emph{Contradicci\'on}. Luego $\SPACE(n)$
\textbf{no} es cerrada.
\end{proof}

\begin{porque}
Re\'une las dos piezas: $B$ es ``$A$ disfrazado de f\'acil'' por el relleno, y la
reducci\'on $f$ es solo ``poner el disfraz'', que es trivial ($\log$ espacio). Si la
clase respetara las reducciones, el disfraz har\'ia entrar a $A$ en $\SPACE(n)$ ---
pero sabemos que $A$ es genuinamente dif\'icil. La \'unica salida l\'ogica es que
$\SPACE(n)$ \emph{no} respeta las reducciones.
\end{porque}

\subsection*{Conclusi\'on: $\NP\neq\SPACE(n)$}

\begin{proof}
Tenemos:
\begin{itemize}[nosep]
  \item $\NP$ \emph{s\'i} es cerrada para $\log$-reducciones (Paso 1);
  \item $\SPACE(n)$ \emph{no} lo es (Paso 2).
\end{itemize}
El cierre para $\log$-reducciones es una propiedad que una clase tiene o no tiene; dos
clases iguales coincidir\'ian en ella. Como $\NP$ la tiene y $\SPACE(n)$ no,
forzosamente
\[
\boxed{\ \NP\neq\SPACE(n).\ }
\]
\end{proof}

\begin{nota}[recapitulaci\'on del esquema, por si cae en examen]
El patr\'on completo es: (1) elige una propiedad ``$P$ es cerrada para cierto tipo de
reducci\'on''; (2) demuestra que una clase la cumple (aqu\'i, v\'ia certificado +
``$\log$-espacio $\Rightarrow$ tiempo polin\'omico''); (3) demuestra que la otra no,
\emph{exhibiendo} un contraejemplo $A\reduce B$ con $B$ dentro y $A$ fuera, donde la
separaci\'on $A\notin$clase viene de un teorema de jerarqu\'ia y la reducci\'on es un
\emph{padding} trivial; (4) concluye que las clases difieren. Los \'unicos
ingredientes ``pesados'' son el teorema de jerarqu\'ia (diagonalizaci\'on) y la idea
de padding; el resto es encadenar definiciones.
\end{nota}

\vspace{6mm}\hrule\vspace{3mm}
\begin{center}\small
Fin de la Relaci\'on 4 (versi\'on comentada).\\[1mm]
Documento de estudio. Resultados est\'andar en teor\'ia de la complejidad
(Sipser; Papadimitriou; Garey--Johnson); las demostraciones siguen los esquemas
habituales de esos textos y la terminolog\'ia de la asignatura.
\end{center}

\end{document}
